\documentclass[a4paper,11pt,twoside]{report}

% Japanese support with LuaLaTeX
\usepackage{luatexja}
\usepackage{luatexja-fontspec}
\setmainjfont{Hiragino Mincho ProN}
\setsansjfont{Hiragino Kaku Gothic ProN}

% Essential packages
\usepackage{textcomp}
\usepackage[margin=2.5cm]{geometry}
\setlength{\headheight}{14pt}
\usepackage{graphicx}
\usepackage{xcolor}
\usepackage{titlesec}
\usepackage{fancyhdr}
\usepackage{booktabs}
\usepackage{longtable}
\usepackage{array}
\usepackage{multirow}
\usepackage{hyperref}
\usepackage{enumitem}
\usepackage{tikz}
\usepackage{tcolorbox}
\usepackage{fontspec}
\usepackage{parskip}

% Colors
\definecolor{miyabiprimary}{RGB}{16,185,129}
\definecolor{miyabisecondary}{RGB}{59,130,246}
\definecolor{miyabiaccent}{RGB}{245,158,11}
\definecolor{miyabidark}{RGB}{31,41,55}
\definecolor{miyabilight}{RGB}{249,250,251}

% Hyperref setup
\hypersetup{
    colorlinks=true,
    linkcolor=miyabiprimary,
    urlcolor=miyabisecondary,
    citecolor=miyabiaccent
}

% Header/Footer
\pagestyle{fancy}
\fancyhf{}
\fancyhead[LE,RO]{\leftmark}
\fancyhead[RE,LO]{Miyabi Plugin Marketplace}
\fancyfoot[C]{\thepage}
\renewcommand{\headrulewidth}{0.4pt}
\renewcommand{\footrulewidth}{0.4pt}

% Chapter/Section formatting
\titleformat{\chapter}[display]
{\normalfont\huge\bfseries\color{miyabidark}}
{\chaptertitlename\ \thechapter}{20pt}{\Huge}

\titleformat{\section}
{\normalfont\Large\bfseries\color{miyabiprimary}}
{\thesection}{1em}{}

\titleformat{\subsection}
{\normalfont\large\bfseries\color{miyabisecondary}}
{\thesubsection}{1em}{}

% Box styles
\tcbset{
    colback=miyabilight,
    colframe=miyabiprimary,
    fonttitle=\bfseries,
    arc=3mm,
    boxrule=0.5mm
}

\begin{document}

% =============================================================================
% TITLE PAGE
% =============================================================================
\begin{titlepage}
\centering
\vspace*{2cm}

{\Huge\bfseries\color{miyabidark} Miyabi Plugin Marketplace}\\[0.5cm]
{\Large\color{miyabiprimary} 12 Phase Business Strategy Report}\\[2cm]

{\LARGE 包括的ビジネス戦略書}\\[0.3cm]
{\large Year 1-3 完全実行計画}\\[2cm]

\begin{tikzpicture}
\fill[miyabiprimary!20] (0,0) circle (3cm);
\fill[miyabisecondary!30] (0,0) circle (2cm);
\fill[miyabiaccent!40] (0,0) circle (1cm);
\end{tikzpicture}

\vfill

{\large\textbf{作成日}: 2025年11月29日}\\[0.3cm]
{\large\textbf{バージョン}: 1.0.0}\\[0.3cm]
{\large\textbf{作成}: Miyabi AI Agent Orchestra}\\[1cm]

\begin{tabular}{|l|l|}
\hline
\textbf{Year 3 目標} & \textbf{達成値} \\
\hline
ARR & ¥ 1,896M \\
Pro Users & 10,000名 \\
Enterprise & 200社 \\
LTV/CAC & 18.5x \\
NPS & 65+ \\
\hline
\end{tabular}

\end{titlepage}

% =============================================================================
% TABLE OF CONTENTS
% =============================================================================
\tableofcontents
\newpage

% =============================================================================
% EXECUTIVE SUMMARY
% =============================================================================
\chapter{エグゼクティブサマリー}

\section{製品概要}

\textbf{Miyabi Plugin Marketplace}は、AI開発者向けのプラグインエコシステムです。

\begin{tcolorbox}[title=ワンライナー]
「1つのプラグインで、全てのAI IDEに対応。開発時間を10倍削減し、ROI 23倍を実現。」
\end{tcolorbox}

\subsection{市場機会}
\begin{itemize}
    \item AI開発ツール市場: \$50B（2025年）、CAGR 35\%
    \item ターゲット: AI開発者 500万人（有料転換率 3-5\%想定）
\end{itemize}

\subsection{収益予測}

\begin{table}[h]
\centering
\begin{tabular}{lrrr}
\toprule
\textbf{Year} & \textbf{ARR} & \textbf{YoY成長率} & \textbf{営業利益率} \\
\midrule
Year 1 & ¥ 71.28M & N/A & -112\% \\
Year 2 & ¥ 415.2M & +482\% & +10.5\% \\
Year 3 & ¥ 1,896M & +356.5\% & +53.4\% \\
\bottomrule
\end{tabular}
\caption{3年収益予測}
\end{table}

\section{3つの差別化ポイント}

\subsection{1. マルチプラットフォーム対応}
\begin{itemize}
    \item Universal AI IDE Plugin Adapter（特許出願中）
    \item 対応IDE: Claude Code, GitHub Copilot, Cursor, VSCode
    \item 開発コスト75\%削減
\end{itemize}

\subsection{2. AIエージェント統合}
\begin{itemize}
    \item 25 AI Agents（9 Coding + 16 Business）
    \item Issue → PR完成: 3時間 → 15分（12倍高速）
    \item 並列実行で効率72\%向上
\end{itemize}

\subsection{3. TCGスタイルの楽しい体験}
\begin{itemize}
    \item プラグインをキャラクター化（ポケモン図鑑スタイル）
    \item ゲーミフィケーション: Exp獲得、レベルアップ
    \item ユーザーエンゲージメント3倍
\end{itemize}

% =============================================================================
% PHASE 1-4: FOUNDATION
% =============================================================================
\chapter{Phase 1-4: ファンデーション}

\section{Phase 1: 自己分析}

\begin{tcolorbox}[title=SelfAnalysisAgent による分析]
\textbf{強み}:
\begin{itemize}
    \item Universal Adapter技術（特許出願中）
    \item 25 AI Agents統合
    \item 500+ Issue処理実績
\end{itemize}

\textbf{機会}:
\begin{itemize}
    \item AI開発ツール市場: 年率35\%成長
    \item リモートワーク普及によるプラグイン需要増
\end{itemize}
\end{tcolorbox}

\section{Phase 2: 市場調査}

\subsection{TAM/SAM/SOM分析}

\begin{table}[h]
\centering
\begin{tabular}{lrl}
\toprule
\textbf{指標} & \textbf{規模} & \textbf{定義} \\
\midrule
TAM & \$50B & AI開発ツール市場全体 \\
SAM & \$5B & プラグイン/拡張機能市場 \\
SOM & \$100M & Year 3時点の獲得可能市場 \\
\bottomrule
\end{tabular}
\caption{市場規模分析}
\end{table}

\section{Phase 3: ペルソナ設計}

\subsection{ペルソナ1: たかしさん（28歳・フリーランス）}
\begin{itemize}
    \item 課題: タスク管理苦手、コードレビュー時間
    \item 欲しいもの: 全IDE対応、簡単操作
    \item 利用プラン: Pro（¥ 9,800/月）
\end{itemize}

\subsection{ペルソナ2: さゆりさん（35歳・スタートアップCTO）}
\begin{itemize}
    \item 課題: チーム開発速度、品質バラつき
    \item 欲しいもの: ROI測定可能なツール
    \item 利用プラン: Enterprise（¥ 300,000/月〜）
\end{itemize}

\section{Phase 4: プロダクトコンセプト}

\subsection{ビジネスモデルキャンバス概要}

\textbf{Value Proposition}:
\begin{itemize}
    \item 開発時間10倍削減
    \item マルチプラットフォーム対応
    \item AI Agent統合
    \item TCGスタイル体験
\end{itemize}

\textbf{Revenue Streams}:
\begin{itemize}
    \item Subscription: 85\%
    \item Transaction Fee: 10\%
    \item Enterprise Services: 5\%
\end{itemize}

% =============================================================================
% PHASE 5-6: PRODUCT
% =============================================================================
\chapter{Phase 5-6: プロダクト設計}

\section{Phase 5: サービス詳細設計}

\subsection{価格設定}

\begin{table}[h]
\centering
\begin{tabular}{llr}
\toprule
\textbf{プラン} & \textbf{機能} & \textbf{価格} \\
\midrule
Free & プラグイン5個まで & ¥ 0 \\
Pro & 全機能 + 25 Agents & ¥ 9,800/月 \\
Enterprise & 専任CSM + SLA & ¥ 300,000/月〜 \\
\bottomrule
\end{tabular}
\caption{価格プラン}
\end{table}

\subsection{Unit Economics}

\begin{table}[h]
\centering
\begin{tabular}{lrrr}
\toprule
\textbf{指標} & \textbf{Year 1} & \textbf{Year 2} & \textbf{Year 3} \\
\midrule
LTV (Pro) & ¥ 326,666 & ¥ 326,666 & ¥ 326,666 \\
CAC & ¥ 145,161 & ¥ 43,902 & ¥ 17,647 \\
LTV/CAC & 2.25x & 7.44x & 18.5x \\
Payback & 17.4ヶ月 & 5.3ヶ月 & 2.1ヶ月 \\
\bottomrule
\end{tabular}
\caption{Unit Economics推移}
\end{table}

\section{Phase 6: コンテンツ制作}

\subsection{コンテンツカレンダー（6ヶ月）}

\begin{itemize}
    \item YouTube: 週2本（72動画）
    \item ブログ: 週2本（48記事）
    \item note.com: 週3本（週末配信）
    \item Udemy: 3コース
\end{itemize}

% =============================================================================
% PHASE 7-8: MARKETING
% =============================================================================
\chapter{Phase 7-8: マーケティング戦略}

\section{Phase 7: ファネル設計}

\subsection{6ステージファネル}

\begin{table}[h]
\centering
\begin{tabular}{lrrl}
\toprule
\textbf{ステージ} & \textbf{人数} & \textbf{転換率} & \textbf{主要施策} \\
\midrule
認知 & 100,000 & - & SEO, X, Product Hunt \\
興味 & 10,000 & 10\% & LP最適化, 動画デモ \\
検討 & 3,000 & 30\% & 無料プラン, オンボーディング \\
購入 & 300 & 10\% & 限定割引, ROI計算ツール \\
継続 & 240 & 80\% & CS強化, コミュニティ \\
推奨 & 90 & 30\% & 紹介報酬, レアキャラ \\
\bottomrule
\end{tabular}
\caption{ファネル設計（Month 6目標）}
\end{table}

\section{Phase 8: SNS戦略}

\subsection{プラットフォーム別目標（Month 6）}

\begin{table}[h]
\centering
\begin{tabular}{lr}
\toprule
\textbf{プラットフォーム} & \textbf{目標} \\
\midrule
X (Twitter) & 5,000フォロワー \\
YouTube & 10,000登録者 \\
GitHub Stars & 2,500 \\
Discord & 1,000メンバー \\
dev.to & 1,500フォロワー \\
LinkedIn & 2,000フォロワー \\
\bottomrule
\end{tabular}
\caption{SNS目標}
\end{table}

\subsection{コンテンツミックス}
\begin{itemize}
    \item 教育・ノウハウ: 40\%
    \item トレンド参加: 25\%
    \item ブランドストーリー: 20\%
    \item プロモーション: 10\%
    \item コミュニティ: 5\%
\end{itemize}

% =============================================================================
% PHASE 9-10: SALES
% =============================================================================
\chapter{Phase 9-10: セールス戦略}

\section{Phase 9: マーケティング施策}

\subsection{7P分析}

\begin{table}[h]
\centering
\begin{tabular}{ll}
\toprule
\textbf{要素} & \textbf{戦略} \\
\midrule
Product & マルチプラットフォーム対応プラグイン \\
Price & Value-Based Pricing (¥ 9,800/月) \\
Place & オンラインMarketplace \\
Promotion & PLG + コンテンツマーケティング \\
People & 専任CSM配置 \\
Process & 自動化95\% (Free/Pro) \\
Physical & ダッシュボード・分析ツール \\
\bottomrule
\end{tabular}
\caption{7P Marketing Mix}
\end{table}

\section{Phase 10: セールス戦略}

\subsection{セグメント別セールスモーション}

\begin{table}[h]
\centering
\begin{tabular}{llrl}
\toprule
\textbf{セグメント} & \textbf{モーション} & \textbf{サイクル} & \textbf{自動化率} \\
\midrule
SMB (1-10名) & Self-Serve PLG & 0-7日 & 95\% \\
Mid-Market (11-100名) & Sales-Assisted & 14-30日 & 40\% \\
Enterprise (100+名) & Field Sales & 60-180日 & 10\% \\
\bottomrule
\end{tabular}
\caption{セールスモーション設計}
\end{table}

\subsection{パイプライン転換率}

\begin{table}[h]
\centering
\begin{tabular}{lr}
\toprule
\textbf{ステージ遷移} & \textbf{転換率} \\
\midrule
MQL → SQL & 40\% \\
SQL → Opportunity & 60\% \\
Opportunity → Proposal & 70\% \\
Proposal → Negotiation & 80\% \\
Negotiation → Closed Won & 90\% \\
\midrule
\textbf{MQL → Closed Won} & \textbf{12.1\%} \\
\bottomrule
\end{tabular}
\caption{パイプライン転換率}
\end{table}

% =============================================================================
% PHASE 11-12: OPERATIONS
% =============================================================================
\chapter{Phase 11-12: オペレーション}

\section{Phase 11: CRM・顧客管理}

\subsection{CS組織構成（Year 3）}

\begin{table}[h]
\centering
\begin{tabular}{lrl}
\toprule
\textbf{役職} & \textbf{人数} & \textbf{担当} \\
\midrule
VP of CS & 1名 & 部門統括 \\
CSM Lead & 3名 & Enterprise/Pro/Self-Service \\
CSM & 10名 & Enterprise 4名, Pro 6名 \\
Support Engineer & 2名 & 技術サポート \\
Community Manager & 1名 & Discord運営 \\
\midrule
\textbf{合計} & \textbf{17名} & \\
\bottomrule
\end{tabular}
\caption{CS組織構成}
\end{table}

\subsection{Health Score設計}

\begin{table}[h]
\centering
\begin{tabular}{lrl}
\toprule
\textbf{要素} & \textbf{ウェイト} & \textbf{計測項目} \\
\midrule
利用頻度 & 20\% & DAU, 実行回数 \\
エンゲージメント & 30\% & 機能採用率, ログイン頻度 \\
サポート満足度 & 20\% & CSAT, 応答時間 \\
支払い状況 & 15\% & 遅延なし, 更新率 \\
NPS & 15\% & 四半期調査 \\
\bottomrule
\end{tabular}
\caption{Health Score構成}
\end{table}

\subsection{リスクレベル対応}

\begin{table}[h]
\centering
\begin{tabular}{lll}
\toprule
\textbf{スコア} & \textbf{リスク} & \textbf{対応} \\
\midrule
0-40 & High Risk & 即座に介入（24時間以内） \\
41-60 & Medium Risk & CSMコンタクト（48時間以内） \\
61-80 & Healthy & モニタリング継続 \\
81-100 & Champion & Advocate育成プログラム \\
\bottomrule
\end{tabular}
\caption{リスクレベル対応}
\end{table}

\section{Phase 12: データ分析・PDCA}

\subsection{データアーキテクチャ}

\textbf{データソース（8種類）}:
\begin{itemize}
    \item GA4（Webトラフィック）
    \item 広告データ（Google/Meta/LinkedIn）
    \item CRMデータ（HubSpot）
    \item 売上データ（Stripe）
    \item SNSデータ（X/Discord/YouTube）
    \item プロダクトデータ（Mixpanel）
    \item サポートデータ（Intercom）
    \item NPS調査（Delighted）
\end{itemize}

\textbf{インフラストラクチャ}:
\begin{itemize}
    \item DWH: BigQuery
    \item BI: Looker Studio
    \item Product Analytics: Mixpanel
    \item Alerting: PagerDuty
\end{itemize}

\subsection{North Star Metric}

\begin{tcolorbox}[title=North Star Metric]
\textbf{Weekly Active Plugins (WAP)}

週次アクティブプラグイン数を主要指標として採用。

目標:
\begin{itemize}
    \item Month 3: 5,000 WAP
    \item Month 6: 50,000 WAP
    \item Year 1: 500,000 WAP
\end{itemize}

理由: プラグイン実行数 = ユーザー価値実現。有料転換率と高相関（R² = 0.85）
\end{tcolorbox}

% =============================================================================
% FINANCIAL PROJECTIONS
% =============================================================================
\chapter{財務計画}

\section{3年収益予測}

\begin{table}[h]
\centering
\begin{tabular}{lrrr}
\toprule
\textbf{指標} & \textbf{Year 1} & \textbf{Year 2} & \textbf{Year 3} \\
\midrule
Free Users & 10,000 & 50,000 & 200,000 \\
Pro Users & 300 & 2,000 & 10,000 \\
Enterprise & 10社 & 50社 & 200社 \\
\midrule
ARR & ¥ 71.28M & ¥ 415.2M & ¥ 1,896M \\
YoY成長率 & N/A & +482\% & +356.5\% \\
\midrule
売上原価 & ¥ 10.69M & ¥ 62.28M & ¥ 284.4M \\
粗利益 & ¥ 60.59M & ¥ 352.92M & ¥ 1,611.6M \\
\midrule
販売管理費 & ¥ 140.71M & ¥ 309.4M & ¥ 599.6M \\
営業利益 & -¥ 80.12M & +¥ 43.52M & +¥ 1,012M \\
営業利益率 & -112\% & +10.5\% & +53.4\% \\
\bottomrule
\end{tabular}
\caption{P/L予測}
\end{table}

\section{資金調達計画}

\subsection{Year 1: シードラウンド}
\begin{itemize}
    \item 調達額: ¥ 200,000,000
    \item Pre-money Valuation: ¥ 800,000,000
    \item 用途: プロダクト開発40\%, マーケティング30\%, 人材採用20\%, 運転資金10\%
    \item Runway: 12-18ヶ月
\end{itemize}

\subsection{Year 2: シリーズA}
\begin{itemize}
    \item 調達額: ¥ 500,000,000
    \item Pre-money Valuation: ¥ 4,000,000,000
    \item 用途: マーケティング強化50\%, 人材採用30\%, グローバル展開20\%
    \item Runway: 24ヶ月
\end{itemize}

\subsection{Year 3: 黒字化}
\begin{itemize}
    \item 営業CF: +¥ 1,012,000,000
    \item 追加調達: 不要
    \item 戦略: 利益を成長投資に再配分
\end{itemize}

\section{Rule of 40}

\begin{table}[h]
\centering
\begin{tabular}{lrr}
\toprule
\textbf{Year} & \textbf{成長率 + 利益率} & \textbf{評価} \\
\midrule
Year 1 & N/A & - \\
Year 2 & 482\% + (-5\%) = 477\% & Excellent \\
Year 3 & 356.5\% + 15\% = 371.5\% & Excellent \\
\bottomrule
\end{tabular}
\caption{Rule of 40分析}
\end{table}

% =============================================================================
% ROADMAP
% =============================================================================
\chapter{6ヶ月ロードマップ}

\section{Month 1-2: MVP構築}

\begin{itemize}
    \item プラグインアップロード機能
    \item MCP v2.0統合
    \item Claude Code + Copilot対応
    \item Stripe決済統合
\end{itemize}

\textbf{KPI}:
\begin{itemize}
    \item betaテスター: 100名
    \item プラグイン: 10個
    \item ダウンロード: 500
\end{itemize}

\section{Month 3-4: パブリックローンチ}

\begin{itemize}
    \item Cursor + VSCode対応
    \item プラグイン検索強化
    \item レビューシステム
    \item SDK公開
\end{itemize}

\textbf{マーケティング}:
\begin{itemize}
    \item Product Hunt
    \item Hacker News
    \item X広告開始
\end{itemize}

\textbf{KPI}:
\begin{itemize}
    \item ユーザー: 1,000名
    \item プラグイン: 50個
    \item Pro転換: 3\%（30名）
    \item MRR: ¥ 294,000
\end{itemize}

\section{Month 5-6: エコシステム拡大}

\begin{itemize}
    \item プラグイン推薦AI（GPT-4o）
    \item カスタムAgent作成UI
    \item チームワークスペース
    \item Enterpriseプラン開始
\end{itemize}

\textbf{パートナーシップ}:
\begin{itemize}
    \item Anthropic公式パートナー申請
    \item Microsoft連携協議
    \item Anysphere連携開始
\end{itemize}

\textbf{KPI}:
\begin{itemize}
    \item ユーザー: 5,000名
    \item プラグイン: 200個
    \item Pro転換: 4\%（200名）
    \item Enterprise: 3社
    \item MRR: ¥ 2,860,000
\end{itemize}

% =============================================================================
% CONCLUSION
% =============================================================================
\chapter{まとめ}

\section{投資ハイライト}

\begin{tcolorbox}[title=なぜMiyabiに投資すべきか]
\textbf{1. 巨大市場と高成長}
\begin{itemize}
    \item 市場規模: \$50B（2025年）、CAGR 35\%
    \item TAM: AI開発者 500万人
\end{itemize}

\textbf{2. 圧倒的な差別化}
\begin{itemize}
    \item 特許: Universal Adapter（出願中）
    \item 技術: 25 AI Agents統合
    \item UX: TCGスタイル（エンゲージメント3倍）
\end{itemize}

\textbf{3. 強力なUnit Economics}
\begin{itemize}
    \item LTV/CAC: 18.5x（Year 3）
    \item Payback: 2.1ヶ月（Year 3）
    \item Gross Margin: 85\%
\end{itemize}

\textbf{4. 急成長と早期黒字化}
\begin{itemize}
    \item 黒字化: Year 2 Q4
    \item 営業利益率: +53.4\%（Year 3）
    \item Rule of 40: 371.5\%（Year 3）
\end{itemize}
\end{tcolorbox}

\section{Next Steps}

\begin{enumerate}
    \item \textbf{MVP構築開始}（Month 1-2）
    \item \textbf{シード資金調達}（¥ 200M）
    \item \textbf{betaテスター募集}（100名）
    \item \textbf{パブリックローンチ}（Month 3）
\end{enumerate}

\vspace{1cm}

\begin{center}
\Large\bfseries
このコンセプトなら、市場で勝てる。\\
投資家を魅了できる。
\end{center}

\vfill

\begin{center}
\textbf{Miyabi Plugin Marketplace}\\
contact@miyabi-marketplace.com\\
https://miyabi-world.com\\
GitHub: github.com/ShunsukeHayashi/miyabi-private
\end{center}

\end{document}
